% ── PROJECT 01 ──────────────────────────────────────────────────────────────
% DOMAIN: Healthcare AI · Clinical Machine Learning · Predictive Modeling · Published · Surgical Endoscopy 2025
% ATS KEYWORDS: healthcare AI, OR scheduling optimization, surgical duration prediction, perioperative analytics, clinical decision support, machine learning, deep learning, neural network, ANN, XGBoost, random forest, gradient boosting, SVR, predictive modeling, feature engineering, cross-validation, bias analysis, TRIPOD-AI, Python, scikit-learn, Keras, TensorFlow, Bayesian hyperparameter optimization, 10-fold cross-validation, hospital operations research
% TECHNICAL DETAILS: Retrospective cohort study using 17,246 elective surgical records (16,159 unique patients, 2015–2020) from 3 academic tertiary hospitals in London, Canada (Cerner EMR). Benchmarked 8 ML models — Linear Regression, Ridge, Lasso, SVR, Random Forest, GBM, XGBoost, ANN — to predict OR case duration from 20 pre-operative structured features (procedure type, ASA score, age, BMI, surgeon identity, ICD code, day-of-week). Best model: ANN, MAE=31.8 min (vs. 35.3 min surgeon baseline), MAPE=26% (vs. 34%), R²=0.78, systematic bias=−0.37 min (p=0.34, statistically zero). Surgeon estimates showed −18.52 min underestimation (p<0.001). Validated with 10-fold cross-validation, Bayesian hyperparameter tuning, TRIPOD-AI standards. Tech stack: Python, scikit-learn, Keras/TensorFlow, XGBoost, Bayesian Optimization.
% NARRATIVE (from portfolio): Benchmarked 8 ML models on 17,246 real elective surgical records from 3 academic hospitals to predict OR case duration from 20 structured pre-operative features. The neural network (ANN) achieved the best performance — MAE=31.8 min, MAPE=26%, R²=0.78 — and was the only model with statistically zero scheduling bias (−0.37 min, p=0.34), while surgeon estimates showed a systematic −18.52 min underestimation (p<0.001). Validated using 10-fold cross-validation, Bayesian hyperparameter tuning, and TRIPOD-AI reporting standards. Tech stack: Python, scikit-learn, XGBoost, Keras/TensorFlow, Bayesian Optimization.
% HTML TAGS: Neural Network, ANN, XGBoost, Random Forest, GBM, scikit-learn, Bayesian Optimization, TRIPOD-AI, Healthcare AI, Predictive Modeling, Feature Engineering, Python
% \cvprojectnogit{ML Benchmark: Surgical OR Duration Prediction}{}{}{https://mnoorchenar.github.io/projects/Projects-Files/001-surgical-duration.html}{Neural network trained on 17,246 surgical records across 3 hospitals predicts OR duration with near-zero systematic bias ($-$0.37 min, p=0.34), outperforming surgeon estimates by 18+ min; benchmarked against 7 other ML models with 10-fold CV and Bayesian hyperparameter tuning. Published in \textit{Surgical Endoscopy} (2025).}{\ptag{Healthcare AI}\ptag{Neural Network}\ptag{XGBoost}\ptag{scikit-learn}\ptag{Bayesian Opt.}\ptag{TRIPOD-AI}\ptag{Python}}

% ── PROJECT 02 ──────────────────────────────────────────────────────────────
% DOMAIN: Natural Language Processing · Model Evaluation · SSRN Preprint 2025
% ATS KEYWORDS: NLP, clinical NLP, transformer embeddings, Sentence-BERT, ClinicalBERT, HuggingFace, text encoding, TF-IDF, label encoding, count vectorization, multimodal learning, feature fusion, surgical duration prediction, healthcare AI, perioperative analytics, XGBoost, scikit-learn, PyTorch, SHAP, explainability, regression, predictive modeling, benchmarking, NLP evaluation, clinical text mining
% TECHNICAL DETAILS: Systematic benchmark of 5 text-encoding strategies — Label Encoding, Count Vectorization, TF-IDF, ClinicalBERT (MIMIC-III fine-tuned), Sentence-BERT — applied to free-text operative notes from 180,366 real elective surgical cases across 3 tertiary academic hospitals in London, Canada (2015–2020). Each strategy fused with structured features (age, BMI, ASA score, sex, procedure type) and evaluated against multiple ML models (Linear Regression, XGBoost, Neural Networks) using MAE, SMAPE, R² with 10-fold cross-validation. Best: Sentence-BERT MAE=27.6 min, SMAPE=23.0%, R²=0.77, reducing prediction error by up to 16% over traditional encodings. ClinicalBERT showed statistically significant gains over structured-only baseline (p<0.01). SHAP used for feature attribution. Preprint: SSRN DOI 10.2139/ssrn.5489578.
% NARRATIVE (from portfolio): Conducted a large-scale NLP benchmarking study comparing five encoding strategies — Label Encoding, Count Vectorization, TF-IDF, ClinicalBERT, and Sentence-BERT — on 180,366 surgical operative notes for duration prediction, quantifying the performance uplift from domain-specific and sentence-level pre-training. Each strategy evaluated on identical downstream tasks using standardized Scikit-learn pipelines for fair, reproducible comparison. Results confirm Sentence-BERT outperforms all other strategies by capturing full sentence-level meaning; ClinicalBERT also showed statistically significant gains over traditional encodings, validating domain-adapted pre-training for production healthcare NLP systems.
% HTML TAGS: NLP, Sentence-BERT, ClinicalBERT, TF-IDF, HuggingFace, Transformers, scikit-learn, PyTorch, XGBoost, SHAP, Clinical Text Mining, Benchmarking, Python
% \cvprojectnogit{Clinical NLP Benchmark: Sentence-BERT, ClinicalBERT, TF-IDF}{}{}{https://mnoorchenar.github.io/projects/Projects-Files/002-nlp-benchmark.html}{Benchmarked 5 text-encoding strategies (Label Encoding to Sentence-BERT) on 180,366 surgical operative notes fused with structured clinical features; Sentence-BERT achieved MAE=27.6 min --- a 16\% error reduction over traditional encodings; SHAP used for feature attribution across all model--encoding combinations. Preprint on SSRN (2025).}{\ptag{NLP}\ptag{Sentence-BERT}\ptag{ClinicalBERT}\ptag{HuggingFace}\ptag{PyTorch}\ptag{XGBoost}\ptag{SHAP}}


% ── PROJECT 03 ──────────────────────────────────────────────────────────────
% DOMAIN: Explainable AI · Time Series · Anomaly Detection · Energy & Buildings 2025
% ATS KEYWORDS: explainable AI, XAI, SHAP, LIME, context-aware explainability, anomaly detection, smart meter, energy analytics, time series, deep learning, LSTM, GRU, BiLSTM, BiGRU, 1D-CNN, TCN, DCNN, WaveNet, Temporal Fusion Transformer, Time Series Transformer, feature attribution, interpretability, Python, cosine similarity, Random Forest, IQR thresholding, Bayesian hyperparameter optimization, NSERC, London Hydro, min-max scaling, 48-hour sliding window
% TECHNICAL DETAILS: Published Energy & Buildings (Vol. 328, 2025, DOI: 10.1016/j.enbuild.2024.115177). Context-aware SHAP baseline selection replacing random background with cosine-similarity-matched historical windows, weighted by Random Forest feature importances. Evaluated across 5 building types (residential, manufacturing, medical clinic, retail, office) using datasets spanning 2002–2017, 10 DL architectures (LSTM, GRU, BiLSTM, BiGRU, 1D-CNN, TCN, DCNN, WaveNet, TFT, TST), 5 XAI techniques (Kernel SHAP, Partition SHAP, Sampling SHAP, LIME, Permutation Importance), 250+ model–dataset–XAI combinations. Results: 38% avg reduction in explanation variability, max 80.3% (manufacturing, Partition SHAP, p<0.001). IQR-based thresholding on prediction residuals with 48-hour sliding window inputs and min-max scaling. Bayesian hyperparameter tuning on all 10 models. Funded by NSERC and London Hydro.
% NARRATIVE (from portfolio): Developed a context-aware XAI framework for deep learning anomaly detection in smart-meter energy data, replacing random SHAP baselines with cosine-similarity-matched historical windows weighted by Random Forest feature importances. Evaluated across 5 building types, 10 deep learning architectures (LSTM, GRU, BiLSTM, BiGRU, 1D-CNN, TCN, DCNN, WaveNet, TFT, TST), and 5 XAI methods over 250+ combinations. Achieved 38% average reduction in explanation variability (max 80.3%, manufacturing, Partition SHAP, p<0.001). Pipeline uses IQR-based anomaly thresholding with 48-hour sliding window inputs, min-max scaling, and Bayesian hyperparameter tuning.
% HTML TAGS: Explainable AI, XAI, SHAP, LIME, LSTM, GRU, TCN, WaveNet, Deep Learning, Anomaly Detection, Time Series, Energy Analytics, Cosine Similarity, Random Forest, Python
% \cvprojectnogit{Explainable AI for Deep Learning Anomaly Detection}{}{}{https://mnoorchenar.github.io/projects/Projects-Files/003-energy-anomaly-xai.html}{Context-aware SHAP framework for deep learning anomaly detection in smart-meter energy data; replaces random SHAP baselines with cosine-similarity-matched historical windows, reducing explanation variability by 38\% on average (up to 80.3\%) across 5 building types and 10 architectures (LSTM, GRU, BiLSTM, BiGRU, 1D-CNN, TCN, DCNN, WaveNet, TFT, TST). Published in \textit{Energy \& Buildings} (2025).}{\ptag{XAI}\ptag{SHAP}\ptag{LSTM}\ptag{TCN}\ptag{WaveNet}\ptag{Anomaly Detection}\ptag{Python}}


% ── PROJECT 04 ──────────────────────────────────────────────────────────────
% DOMAIN: Biostatistics · Genomics · Genetic Risk Prediction · J. Biostatistics & Epidemiology 2022
% ATS KEYWORDS: biostatistics, genomics, GWAS, candidate gene, MultiPhen, GATES, SNP, linkage disequilibrium, ordinal regression, sarcopenia, muscle loss, aging, IL10, interleukin, inflammation, population genetics, non-European cohort, Middle Eastern, BEH cohort, Tehran University, R, PLINK, multivariate statistics, public health, epidemiology, Extended Simes procedure, HWE filtering, MAF filtering, joint regression, gene-level aggregation
% TECHNICAL DETAILS: Published J. Biostatistics & Epidemiology, 8(2): 195–207, 2022. BEH cohort, Iran — 2,772 participants aged 60+, 663,377 genome-wide SNPs. MultiPhen joint ordinal regression of genotype on both SMI and Handgrip Strength simultaneously (up to 2× statistical power vs. single-trait GWAS). GATES gene-level confirmation aggregating 27 independent SNPs (LD filtering: r²≤0.4, MAF>0.01, HWE p>0.05) within ±50kb of IL10 (chr 1). Results: 3 significant intronic IL10 variants (rs11119603 p=0.00384, rs57461190 p=0.00411, rs3950619 p=0.03641); GATES p=0.046 — first IL10 confirmation as sarcopenia risk gene in an Iranian population. Effect sizes 0.178–0.883. Tech stack: R, MultiPhen, GATES, PLINK.
% NARRATIVE (from portfolio): Applied MultiPhen joint ordinal regression and GATES gene-level aggregation on 663,377 SNPs from 2,772 elderly Iranians (BEH cohort) to identify IL10 as the first genomically validated sarcopenia risk gene in an Iranian population. The joint multivariate approach — testing both skeletal muscle index and handgrip strength simultaneously — provided up to 2× more statistical power than single-trait GWAS. Three intronic IL10 variants confirmed individually (p=0.00384–0.03641); GATES gene-level confirmation at p=0.046. Tech stack: R, MultiPhen, GATES, PLINK.
% HTML TAGS: GWAS, Biostatistics, MultiPhen, GATES, Ordinal Regression, Linkage Disequilibrium, IL10, Candidate-Gene, Population Genetics, Epidemiology, PLINK, R
% \cvprojectnogit{GWAS Identifies IL10 as Sarcopenia Risk Gene}{}{}{https://mnoorchenar.github.io/projects/Projects-Files/004-sarcopenia-genetics.html}{Joint MultiPhen + GATES analysis of 663K SNPs in 2,772 elderly Iranians confirmed IL10 (p=0.046) as the first genomically validated sarcopenia risk gene in an Iranian population; multivariate joint regression of both SMI and handgrip strength yielded up to 2$\times$ more statistical power than single-trait GWAS. Published in \textit{J. Biostatistics \& Epidemiology} (2022).}{\ptag{GWAS}\ptag{MultiPhen}\ptag{GATES}\ptag{Biostatistics}\ptag{R}\ptag{Epidemiology}\ptag{Population Genetics}}


% ── PROJECT 05 ──────────────────────────────────────────────────────────────
% DOMAIN: Metabolomics · Cardiovascular Risk · Biomarker Discovery · Frontiers in Cardiovascular Medicine 2023
% ATS KEYWORDS: metabolomics, FIA-MS/MS, acylcarnitines, amino acids, ASCVD, cardiovascular risk, PCA, Varimax rotation, logistic regression, biomarker discovery, pathway enrichment, KEGG, MetaboAnalyst, BCAA, biostatistics, public health, epidemiology, multivariate statistics, SPSS, Tehran University, Iran, STEPs cohort, Benjamini-Hochberg, FDR, non-European population, KMO test, Bartlett test, ACC/AHA, BMI adjustment
% TECHNICAL DETAILS: Published Frontiers in Cardiovascular Medicine, 10:1161761, 2023. 50 metabolites from 1,102 fasting plasma samples (30 acylcarnitines + 20 amino acids) via FIA-MS/MS (SCIEX API 3200, butanol-HCl derivatisation). Pipeline: ln Z-score standardisation → PCA Varimax rotation (KMO=0.874, Bartlett p<0.001) → 10 orthogonal factors → binary logistic regression (low-risk reference) → MetaboAnalyst v5.0 / KEGG enrichment with Benjamini-Hochberg FDR. 14 significant ASCVD biomarkers: 3 acylcarnitines (C4DC, C8:1, C16OH) + 11 amino acids. Highest OR: Factor 10 (ornithine + citrulline) OR=1.570 (p<0.001). Protective: Factor 9 (glycine, serine, threonine) OR=0.741 (p<0.001). Strongest single-metabolite: C16OH (r=0.279, p<0.001). Software: SPSS v19.0, MetaboAnalyst v5.0. 17 co-authors.
% NARRATIVE (from portfolio): Profiled 50 plasma metabolites (30 acylcarnitines + 20 amino acids) via FIA-MS/MS in 1,102 participants from the STEPs 2016 national survey. Applied PCA with Varimax rotation (KMO=0.874) to reduce 50 correlated metabolites into 10 orthogonal factors, then binary logistic regression across four ACC/AHA ASCVD risk groups with Benjamini-Hochberg FDR correction. KEGG pathway enrichment via MetaboAnalyst v5.0. Identified 14 significant biomarkers (C4DC, C8:1, C16OH + 11 amino acids); Factor 10 (ornithine + citrulline) showed highest risk association (OR=1.57) while Factor 9 (glycine, serine, threonine) emerged as protective (OR=0.74). Tech stack: FIA-MS/MS, SPSS v19.0, MetaboAnalyst v5.0, KEGG.
% HTML TAGS: Metabolomics, FIA-MS/MS, PCA, Logistic Regression, MetaboAnalyst, KEGG, Biomarker Discovery, ASCVD, Cardiovascular Risk, Benjamini-Hochberg, Epidemiology, SPSS
% \cvprojectnogit{Metabolomics Pipeline for ASCVD Biomarker Discovery}{}{}{https://mnoorchenar.github.io/projects/Projects-Files/005-ascvd-metabolomics.html}{PCA (Varimax) and logistic regression on 50 FIA-MS/MS plasma metabolites (1,102 participants) identified 14 significant ASCVD biomarkers --- 3 acylcarnitines (C4DC, C8:1, C16OH) and 11 amino acids; ornithine/citrulline factor showed highest risk association (OR=1.57) while glycine/serine emerged as protective (OR=0.74). Published in \textit{Frontiers in Cardiovascular Medicine} (2023).}{\ptag{Metabolomics}\ptag{FIA-MS/MS}\ptag{PCA}\ptag{Logistic Regression}\ptag{MetaboAnalyst}\ptag{SPSS}\ptag{Epidemiology}}


% ── PROJECT 06 ──────────────────────────────────────────────────────────────
% DOMAIN: Digital Health · Gamification · Randomized Clinical Trial · JMIR 2021
% ATS KEYWORDS: randomized controlled trial, RCT, mHealth, gamification, digital health, patient adherence, CABG, cardiac rehabilitation, teach-back, Android app, biostatistics, ANOVA, Dunnett post-hoc, MMAS, SPSS, STATA, clinical trial, Tehran University, block randomization, health behavior, post-surgical care, Cronbach alpha, Sanaie questionnaire, Fisher exact test
% TECHNICAL DETAILS: Published JMIR, 23(12):e22557, 2021. Three-arm RCT: n=123 CABG patients, block randomization (size 4), n=41 per arm. Arms: (1) gamification — custom Android app "Delban" with diet/medication/movement modules, star reward system (+6/−3), shared social leaderboard; (2) teach-back — 45–60 min session + booklet; (3) usual-care control. Primary outcomes: dietary adherence (Sanaie 30-item, 4-pt Likert), movement adherence (Sanaie 19-item, 5-pt Likert), medication adherence (MMAS, Cronbach α=0.82). One-way ANOVA + Dunnett post-hoc, SPSS v20, STATA v12. Results: diet (F=71.8, p<0.001, Δ=+1.797 vs. control), movement (F=124.5, p<0.001, Δ=+2.013 vs. control), non-overlapping 95% CIs confirming gamification advantage over teach-back. Medication: F=9.66 (p<0.001), no significant difference between gamification and teach-back. Role: Statistical Analysis.
% NARRATIVE (from portfolio): Conducted statistical analysis for a 3-arm RCT (n=123 CABG patients) comparing a custom gamified Android app ("Delban") against teach-back training and usual care over 30 days post-discharge. One-way ANOVA and Dunnett post-hoc tests in SPSS v20 and STATA v12 showed gamification significantly outperformed teach-back on diet (F=71.8, Δ=+1.797, p<0.001) and movement adherence (F=124.5, Δ=+2.013, p<0.001) with non-overlapping 95% CIs. Medication adherence improved in both active arms vs. control (F=9.66, p<0.001) but the two interventions did not differ significantly from each other.
% HTML TAGS: RCT, Gamification, mHealth, Android App, One-Way ANOVA, Dunnett Post-Hoc, MMAS, Block Randomization, Digital Health, SPSS, STATA, Clinical Trial
% \cvprojectnogit{Gamified CABG Recovery vs. Teach-Back: 3-Arm RCT}{}{}{https://mnoorchenar.github.io/projects/Projects-Files/006-cabg-gamification.html}{Three-arm RCT (n=123 CABG patients) comparing a gamified Android app ("Delban") against teach-back and usual care; one-way ANOVA + Dunnett post-hoc confirmed gamification significantly outperformed teach-back on diet (F=71.8, $\Delta$=+1.80) and movement adherence (F=124.5, $\Delta$=+2.01) at 30 days with non-overlapping confidence intervals. Published in \textit{JMIR} (2021).}{\ptag{RCT}\ptag{Gamification}\ptag{mHealth}\ptag{ANOVA}\ptag{Dunnett Post-Hoc}\ptag{SPSS}\ptag{Clinical Trial}}


% ── PROJECT 07 ──────────────────────────────────────────────────────────────
% DOMAIN: Large Language Models · Drug Discovery · Bioinformatics · Generative AI · Live Demo
% ATS KEYWORDS: LLM, large language model, Llama, Mistral, Phi-3, HuggingFace, drug discovery, bioinformatics, protein target identification, UniProt, AlphaFold, PDB, Flask, REST API, NLP, biomedical NLP, prompt engineering, multi-model fallback, API integration, HuggingFace Spaces, drug target, computational biology, Python, JSON, hallucination mitigation, LLM orchestrator, external knowledge integration, structured output generation, biotech AI, UniProt deep search, organism filter, Homo sapiens
% TECHNICAL DETAILS: Flask web application converting plain-language biomedical queries to validated human protein drug targets with 3D structural visualization. Pipeline: (1) 5-model LLM fallback chain (Llama-3.2-3B-Instruct → Mistral-7B-Instruct-v0.3 → Phi-3-mini-4k-instruct → Apriel-5B-Instruct → Llama-3.1-5B-Instruct), each tried with 3 API methods (conversational → chat_completion → text_generation) before cascading; (2) LLMs return structured JSON array of 3–5 human protein targets with UniProt IDs, function summaries, and known drugs; (3) UniProt REST API validation — every LLM-suggested UniProt ID cross-validated with Homo sapiens filter (organism_id:9606), 250M+ proteins searchable; (4) AlphaFold EBI (v4/v3) structure retrieval + RCSB PDB fallback for 3D in-browser rendering; (5) per-protein LLM chat assistant; (6) UniProt deep search fallback with 3-level progressive keyword extraction; (7) curated disease-protein JSON database as last-resort. Deployed HuggingFace Spaces.
% NARRATIVE (from portfolio): Built a Flask web application accepting plain-language biomedical queries and returning validated human protein drug targets with AlphaFold 3D structures rendered in-browser. Core engine is a 5-model LLM fallback chain (Llama-3.2-3B → Mistral-7B → Phi-3-mini → Apriel-5B → Llama-3.1-5B), each tried via 3 API methods before cascading. LLMs return a structured JSON array of 3–5 protein targets; each UniProt ID cross-validated live via UniProt REST API with strict Homo sapiens organism filter. AlphaFold EBI (v4/v3) structures fetched with RCSB PDB fallback. Additional resilience: 3-level UniProt deep search fallback and curated JSON database as last resort. Per-protein LLM chat assistant for follow-up Q&A.
% HTML TAGS: LLM, Flask, HuggingFace, UniProt REST, AlphaFold, RCSB PDB, Llama 3, Mistral, Drug Discovery, Bioinformatics, Generative AI, Python
\cvproject{LLM Drug Target Discovery via AlphaFold Pipeline}{https://github.com/mnoorchenar/Protein-Sequencing-AI-Discovery}{https://huggingface.co/spaces/mnoorchenar/Protein-Sequencing-AI-Discovery}{https://mnoorchenar.github.io/projects/Projects-Files/007-Protein-Sequencing-AI-Discovery.html}{Flask app using a 5-model LLM fallback chain (Llama-3.2-3B, Mistral-7B, Phi-3-mini, Apriel-5B, Llama-3.1-5B), each tried via 3 API methods, to return a JSON array of 3--5 human drug-target proteins per query; each UniProt ID validated live against UniProt REST API (Homo sapiens filter), AlphaFold 3D structures rendered in-browser, and a per-protein LLM chat assistant for follow-up Q\&A. Deployed on HuggingFace Spaces (2025).}{\ptag{LLM}\ptag{Flask}\ptag{HuggingFace}\ptag{UniProt}\ptag{AlphaFold}\ptag{Bioinformatics}\ptag{Drug Discovery}}



% ── PROJECT 08 ──────────────────────────────────────────────────────────────
% DOMAIN: Machine Learning · Financial Risk Management · FinTech · Data Analytics · MLOps · Live Demo
% ATS KEYWORDS: machine learning, ML pipeline, random forest, gradient boosting, logistic regression, scikit-learn, credit risk, fraud detection, market risk, VaR, CVaR, value at risk, expected shortfall, Sharpe ratio, max drawdown, probability of default, PD, LGD, expected loss, loss ratio, combined ratio, underwriting risk, insurance analytics, claims analytics, loan portfolio analysis, concentration risk, risk intelligence, risk dashboard, financial risk management, banking analytics, insurance analytics, FinTech, risk modeling, REST API, microservice, API endpoint, Flask, Plotly, Python, Pandas, NumPy, Docker, HuggingFace Spaces, MLOps, data visualization, interactive dashboard, full-stack, end-to-end ML, synthetic data, feature engineering, historical simulation, production ML, single-file deployment, containerized deployment, train-on-boot, reproducible ML, StandardScaler, sklearn pipeline
% TECHNICAL DETAILS: Full-stack risk intelligence platform built with Flask + Scikit-learn + Plotly, deployed on HuggingFace Spaces via Docker. Seven interactive risk modules: (1) Credit Risk Scorer — Random Forest (100 trees) on 1,200 borrower records, outputs PD, LGD proxy, Expected Loss; (2) Fraud Detection Monitor — Gradient Boosting (100 estimators) on 3,000 transactions, features: amount, hour, foreign flag, velocity, merchant risk (OHE), threshold 0.25; (3) Market Risk Engine — historical simulation on 252 synthetic daily returns, rolling 21-day VaR at 95%/99%, CVaR/Expected Shortfall, Sharpe ratio, max drawdown, $10M synthetic portfolio; (4) Loan Portfolio Concentration Analyser — EL, gross exposure, grade mix, region × purpose default-rate heatmap; (5) Insurance Claims Analytics — 1,000 synthetic policyholders, smoker/BMI/age risk flags, monthly volume trends; (6) Loss & Combined Ratio Monitor — LR vs combined ratio (LR + 25% expense), profitability threshold at LR = 1.0; (7) Underwriting Risk — Logistic Regression (L2) on age, BMI, smoker, children, vehicle age, region (OHE). Three REST API endpoints (/api/credit, /api/fraud, /api/underwriting) for real-time scoring. All models trained at startup on 5,200+ synthetic records; Plotly figures serialised server-side to JSON, rendered client-side. Entire app in a single app.py; containerised via Docker for HuggingFace Spaces (port 7860).
% NARRATIVE (from portfolio): Built a unified banking and insurance risk intelligence platform with seven interactive ML-powered modules in a single Flask application. Three Scikit-learn models (Random Forest for credit risk, Gradient Boosting for fraud detection, Logistic Regression for underwriting) are trained at startup on 5,200+ synthetic records. Market risk is quantified via historical simulation (VaR 95%/99%, CVaR, Sharpe ratio, max drawdown). Three REST API endpoints deliver real-time scoring from browser forms, mirroring production microservice architecture. All Plotly charts are server-side JSON, rendered client-side for responsive interactivity. Entire application — data generation, ML training, 7 pages, REST endpoints, CSS — in one app.py; containerised via Docker and deployed on HuggingFace Spaces.
% HTML TAGS: Machine Learning, Flask, Scikit-learn, Plotly, Random Forest, Gradient Boosting, Credit Risk, Fraud Detection, VaR / CVaR, Risk Analytics, HuggingFace, Docker, REST API, Python, FinTech
\cvproject{RiskSight Pro: Full-Stack Risk Intelligence Platform}{https://github.com/mnoorchenar/risksight-pro}{https://huggingface.co/spaces/mnoorchenar/risksight-pro}{https://mnoorchenar.github.io/projects/Projects-Files/008-risksight-pro.html}{Flask + Scikit-learn platform unifying 7 risk modules (credit scoring, fraud detection, market VaR/CVaR, loan portfolio, claims, underwriting, loss ratio) with 3 live ML models (Random Forest, Gradient Boosting, Logistic Regression) trained on 5,200+ synthetic records; REST API endpoints for real-time scoring; historical simulation engine computing VaR 95\%/99\%, Expected Shortfall, Sharpe ratio, and max drawdown; server-side Plotly JSON charts; single \texttt{app.py} containerised via Docker and deployed on HuggingFace Spaces (2025).}{\ptag{ML}\ptag{Flask}\ptag{Scikit-learn}\ptag{Plotly}\ptag{Credit Risk}\ptag{Fraud Detection}\ptag{VaR/CVaR}\ptag{Risk Analytics}\ptag{Docker}\ptag{FinTech}}


% ── PROJECT 09 ──────────────────────────────────────────────────────────────
% DOMAIN: Generative AI · Deep Learning · Neural Networks · Representation Learning · MLOps · Live Demo
% ATS KEYWORDS: variational autoencoder, VAE, generative model, generative AI, deep learning, neural network, representation learning, latent space, latent representation, dimensionality reduction, encoder, decoder, reparameterisation trick, ELBO, evidence lower bound, KL divergence, binary cross-entropy, BCE, reconstruction loss, manifold learning, unsupervised learning, self-supervised learning, PyTorch, torchvision, MNIST, image generation, image reconstruction, Flask, REST API, base64, Matplotlib, NumPy, Python, HuggingFace Spaces, MLOps, model training, hyperparameter tuning, Adam optimizer, batch training, loss curve, training loop, background thread, non-blocking, interactive ML, data visualization, scatter plot, latent manifold, cluster visualization, generative model exploration, configurable architecture, runtime hyperparameters, deep learning from scratch, neural network architecture, fully connected network, sigmoid output, Gaussian prior, sampling, probabilistic model
% TECHNICAL DETAILS: Interactive browser-based VAE playground built with PyTorch + Flask, deployed on HuggingFace Spaces. Architecture: fully connected VAE, 784 → 400 (configurable hidden dim H) → μ head + log σ² head (configurable latent dim L, default 2) → reparameterise (z = μ + exp(½ log σ²)·ε, ε∼N(0,I)) → 400 → 784 (Sigmoid). Loss: ELBO = BCE (sum) + KL divergence (closed-form Gaussian). Optimizer: Adam, LR=1e-3. Dataset: MNIST, 10K-sample interactive subset (60K available). Training in Python daemon thread; loss curve polled every 600ms via /status endpoint. Seven Flask REST endpoints (/start_training, /status, /latent_space, /reconstruction, /generate, /generate_grid, /architecture) returning base64-encoded PNG responses. Five-tab UI: (1) Live Training Dashboard — configurable epochs, batch size, LR, H, L, live loss curve; (2) Architecture Viewer — dynamic layer diagram + ELBO formula, updates on hyperparameter change; (3) 2-D Latent Manifold — scatter of 10K encoded μ coordinates coloured by digit class; (4) Reconstruction Comparison — 10 random MNIST samples, 2-row side-by-side grid; (5) Latent Navigation — Z₁/Z₂ sliders (−3 to +3) for real-time decoding + 15×15 full manifold grid (225 points). Default hyperparameters: 30 epochs, batch 128, LR 1e-3, H=400, L=2.
% NARRATIVE (from portfolio): Built a fully interactive VAE playground where users configure and train a Variational Autoencoder on MNIST from scratch in the browser, then explore the learned 2-D latent space via scatter plot, side-by-side reconstructions, and slider-driven generation. PyTorch training runs in a background daemon thread; a Flask REST API streams live loss updates every 600ms and returns all visualisations as base64 PNG. Architecture and loss formula update dynamically when hyperparameters change. Seven REST endpoints cover training, latent space visualisation, reconstruction, and generation — all stateless, all base64-encoded. Deployed on HuggingFace Spaces with no local setup required.
% HTML TAGS: Generative AI, VAE, PyTorch, Deep Learning, Representation Learning, ELBO, KL Divergence, Latent Space, MNIST, Flask, REST API, HuggingFace, Neural Networks, Python, Matplotlib
\cvproject{Interactive VAE Playground: Generative AI \& Latent Space Explorer}{https://github.com/mnoorchenar/VAE-Playground}{https://huggingface.co/spaces/mnoorchenar/VAE-Playground}{https://mnoorchenar.github.io/projects/Projects-Files/009-vae-project-page.html}{PyTorch VAE (784→400→L→400→784) trained on MNIST in the browser via a Python daemon thread; ELBO loss (BCE + KL divergence) optimised with Adam; 7 Flask REST endpoints return base64 PNG for a 5-tab UI covering live loss tracking, dynamic architecture viewer, 2-D latent manifold scatter (10K points), side-by-side reconstructions, and real-time latent navigation via Z₁/Z₂ sliders (−3 to +3) with a 15×15 full manifold grid; all hyperparameters (epochs, LR, hidden dim, latent dim) configurable at runtime. Deployed on HuggingFace Spaces (2025).}{\ptag{Generative AI}\ptag{VAE}\ptag{PyTorch}\ptag{Deep Learning}\ptag{Representation Learning}\ptag{Latent Space}\ptag{ELBO}\ptag{Flask}\ptag{HuggingFace}\ptag{Python}}


% ── PROJECT 10 ──────────────────────────────────────────────────────────────
% DOMAIN: Natural Language Processing · Deep Learning · Seq2Seq · Transformer Architecture · MLOps · Live Demo
% ATS KEYWORDS: transformer, transformer architecture, encoder-decoder, seq2seq, sequence-to-sequence, machine translation, neural machine translation, NLP, natural language processing, attention mechanism, multi-head attention, self-attention, cross-attention, scaled dot-product attention, positional encoding, multi-head self-attention, label smoothing, KL divergence, warmup learning rate schedule, inverse-sqrt LR, early stopping, beam search, greedy decoding, BLEU, SacreBLEU, corpus BLEU, perplexity, tokenisation, vocabulary, OOV, out-of-vocabulary, UNK token, special tokens, embedding, feedforward network, layer normalisation, residual connections, dropout, PyTorch, Flask, Server-Sent Events, SSE, Chart.js, Canvas API, Docker, HuggingFace Spaces, Python, NumPy, threading, daemon thread, deep learning from scratch, hyperparameter tuning, configurable architecture, parameter count, attention heatmap, attention visualisation, cross-attention heatmap, English to French, EN→FR, handcrafted dataset, data augmentation, train-val split, checkpoint, best model restoration, MLOps, interactive ML, educational AI, live training, real-time loss streaming
% TECHNICAL DETAILS: 6-step interactive Transformer seq2seq playground built with PyTorch + Flask, deployed on HuggingFace Spaces via Docker. Architecture: vanilla Transformer encoder-decoder; configurable d_model (64–256), num_heads (2/4/8), num_layers (1–4), d_ff (128/256/512), dropout; live parameter count update. Default: d=128, h=4, L=2, d_ff=256, dropout=0.10. Loss: label-smoothed KL divergence (ε=0.05). Optimizer: Adam (β₁=0.9, β₂=0.98), warmup inverse-sqrt LR schedule (100 warmup steps). Dataset: 138 handcrafted EN→FR pairs across 8 vocabulary categories; augmentation slider 1×–16×; 87.5/12.5 train-val split; 15-sentence held-out test set (10 known-vocab, 5 intentional OOV). Training in Python daemon thread; per-epoch SSE JSON events stream loss, PPL, LR, 3 live sample translations to Chart.js live curve. Early stopping with best-checkpoint restoration. Evaluation: SacreBLEU corpus score, greedy vs beam-4 decoding, 10-row side-by-side comparison table. Inference: live OOV token colour-coder (green/red + UNK badge), cross-attention heatmap rendered on HTML Canvas (last decoder layer, averaged across heads). Step-locking enforces Data→Vocab→Model→Train→Evaluate→Infer pipeline order. EN vocab ~95 tokens; FR vocab ~115 tokens; 4 special tokens (PAD, SOS, EOS, UNK). Expected SacreBLEU: 70–95+ at convergence.
% NARRATIVE (from portfolio): Built a 6-step browser-based playground that trains a vanilla Transformer encoder-decoder seq2seq model on a handcrafted 138-pair EN→FR dataset from scratch. Full PyTorch implementation: multi-head attention, positional encoding, label-smoothed KL divergence loss, warmup inverse-sqrt LR. All hyperparameters (d_model, heads, layers, d_ff, dropout) configurable at runtime with live parameter count. Training runs in a daemon thread and streams per-epoch loss, perplexity, LR, and sample translations to the browser via SSE. Post-training: SacreBLEU corpus evaluation with greedy vs beam-4 decoding side-by-side. Inference tab: real-time OOV token flagging, cross-attention heatmap on Canvas (last decoder layer, head-averaged). Step-locking wizard enforces correct conceptual pipeline order. Deployed on HuggingFace Spaces via Docker.
% HTML TAGS: Transformer, NLP, Seq2Seq, PyTorch, Attention Mechanism, BLEU, Beam Search, SSE, Label Smoothing, Flask, Deep Learning, HuggingFace, Docker, Python, Machine Translation
\cvproject{Interactive Transformer Trainer: Seq2Seq \& Attention Visualisation}{https://github.com/mnoorchenar/transformer-edu-viz}{https://huggingface.co/spaces/mnoorchenar/transformer-edu-viz}{https://mnoorchenar.github.io/projects/Projects-Files/010-transformer-edu-viz.html}{PyTorch vanilla Transformer (encoder-decoder seq2seq) trained from scratch on a handcrafted 138-pair EN→FR dataset; configurable d\_model/heads/layers/d\_ff at runtime with live parameter count; label-smoothed KL divergence loss + Adam with warmup inverse-sqrt LR; per-epoch SSE stream drives a Chart.js live loss curve with early stopping; SacreBLEU evaluation with greedy vs beam-4 decoding; cross-attention heatmap rendered on HTML Canvas; real-time OOV token flagging; 6-step step-locked wizard (Data→Vocab→Model→Train→Evaluate→Infer). Deployed on HuggingFace Spaces via Docker (2025).}{\ptag{Transformer}\ptag{NLP}\ptag{Seq2Seq}\ptag{PyTorch}\ptag{Attention}\ptag{BLEU}\ptag{Beam Search}\ptag{SSE}\ptag{Flask}\ptag{Docker}\ptag{HuggingFace}}


% ── PROJECT 11 ──────────────────────────────────────────────────────────────
% DOMAIN: Statistical Analysis · Experimental Design · A/B Testing · Data Science · Decision Intelligence · Live Demo
% ATS KEYWORDS: A/B testing, controlled experiment, statistical experiment, experiment design, experimentation platform, hypothesis testing, null hypothesis, statistical significance, practical significance, p-value, alpha, type I error, type II error, false positive rate, statistical power, power analysis, sample size calculation, minimum detectable effect, MDE, effect size, Cohen's d, two-proportion z-test, Welch's t-test, t-test, z-test, ANOVA, two-way ANOVA, factorial ANOVA, two-factor design, design of experiments, DoE, interaction effect, main effect, replication, sequential testing, interim analysis, peeking problem, alpha spending, O'Brien-Fleming, OBF boundary, Monte Carlo simulation, type I error inflation, multiple testing, SciPy, statsmodels, OLS, NumPy, Plotly, Flask, Flask Blueprint, Jinja2, Docker, HuggingFace Spaces, Python, data-driven decision making, experiment lifecycle, pre-experiment planning, post-experiment analysis, statistical testing, business significance, conversion rate, uplift, treatment and control, traffic split, CSV upload, server-side computation, live computation, debounced input, JSON API, power curve, sample size trade-off, causal inference, experimentation framework, decision framework
% TECHNICAL DETAILS: Full-stack statistical experimentation suite built with Python 3.10 + Flask 3.x, deployed on HuggingFace Spaces via Docker (port 7860). Four self-contained Flask Blueprint modules: (1) Power Calculator — two-proportion z-test power via Fleiss (1981) formula; inputs: baseline conversion, MDE, α, target power, tails; outputs: per-group n, power-vs-MDE curve, sample-size trade-off chart across power targets; (2) A/B Test Analyzer — two-proportion z-test (pooled SE) or Welch's t-test (unequal variance); Cohen's d effect size; 2×2 verdict matrix requiring both statistical significance (p < α) AND user-defined practical significance threshold; supports manual summary stats input and CSV upload; (3) Two-Factor DoE — full two-way factorial ANOVA via statsmodels OLS (Factor A + Factor B + A:B interaction); 2×2 and 3×3 designs; optional replicates per cell; outputs ANOVA table, interaction plot, main-effects chart; (4) Sequential Testing Demo — Monte Carlo simulation (5–50 runs) comparing naive repeated-testing false positive inflation vs O'Brien-Fleming alpha-spending boundary (z_α / √(t/T)). All computation server-side in SciPy + statsmodels + NumPy; results serialised as Plotly graph_objects JSON. Debounced live POST computation (400ms standard, 800ms sequential module). Blueprint architecture: each module self-contained, zero cross-module dependencies.
% NARRATIVE (from portfolio): Built a production-quality statistical experimentation suite covering the full controlled-experiment lifecycle: power analysis (sample sizing and power curves), A/B testing (z-test/t-test with dual statistical + practical significance verdict), two-factor factorial ANOVA with interaction (DoE), and sequential testing (Monte Carlo peeking problem vs O'Brien-Fleming correction). All four modules are self-contained Flask Blueprints with server-side SciPy/statsmodels computation, Plotly JSON charts, debounced live input, and CSV upload support. Dual-gate A/B verdict matrix mirrors industry-standard experimentation decision frameworks. Deployed on HuggingFace Spaces via Docker.
% HTML TAGS: A/B Testing, Statistical Inference, Power Analysis, ANOVA, Design of Experiments, Sequential Testing, SciPy, statsmodels, Plotly, Flask, Docker, HuggingFace, Experimentation, Python, Data Science
\cvproject{ExperimentLab: Statistical Experimentation Suite}{https://github.com/mnoorchenar/experiment-lab}{https://huggingface.co/spaces/mnoorchenar/experiment-lab}{https://mnoorchenar.github.io/projects/Projects-Files/011-experiment-lab.html}{Four-module Flask experimentation platform covering the full controlled-experiment lifecycle: (1) power analysis with MDE/power-curve visualisation; (2) A/B analyzer — two-proportion z-test and Welch's t-test with Cohen's d and a dual-gate 2×2 verdict matrix enforcing both statistical and practical significance; (3) two-way factorial ANOVA with interaction (2×2/3×3, replicate support) via statsmodels OLS; (4) sequential testing — Monte Carlo simulation of peeking-problem false-positive inflation vs O'Brien-Fleming alpha-spending correction; all computation server-side in SciPy/statsmodels; debounced live POST API; Plotly JSON charts; containerised via Docker on HuggingFace Spaces (2024).}{\ptag{A/B Testing}\ptag{Power Analysis}\ptag{ANOVA}\ptag{DoE}\ptag{Sequential Testing}\ptag{SciPy}\ptag{statsmodels}\ptag{Plotly}\ptag{Flask}\ptag{Docker}\ptag{Experimentation}}



% ── PROJECT 12 ──────────────────────────────────────────────────────────────
% DOMAIN: Retrieval-Augmented Generation · Information Retrieval · NLP · LLM Evaluation · MLOps · Live Demo
% ATS KEYWORDS: RAG, retrieval-augmented generation, hybrid search, hybrid retrieval, information retrieval, semantic search, dense retrieval, sparse retrieval, BM25, BM25Okapi, TF-IDF, keyword search, vector search, vector retrieval, FAISS, IndexFlatIP, inner product search, approximate nearest neighbour, ANN, sentence transformer, sentence-transformers, all-MiniLM-L6-v2, embeddings, vector embeddings, dense embeddings, L2 normalisation, cosine similarity, dot product, Reciprocal Rank Fusion, RRF, rank fusion, hybrid ranking, context relevance, faithfulness, diversity, RAG evaluation, LLM evaluation, answer quality, retrieval quality, retrieval benchmark, Flan-T5, seq2seq, generative model, open-source LLM, context distillation, chunking, text chunking, fixed chunking, sentence-window chunking, semantic chunking, chunk overlap, document ingestion, arXiv API, paper ingestion, scientific literature, NLP, natural language processing, question answering, open-domain QA, knowledge retrieval, REST API, JSON API, Flask, Plotly Dash, Docker, HuggingFace Spaces, Python, no API key, open-source, from scratch, end-to-end pipeline, MLOps, hard negative mining, contrastive learning, triplet loss, embedding fine-tuning, domain adaptation, radar chart, evaluation dashboard, token overlap, context window, retrieval pipeline, LLM pipeline, generative AI
% TECHNICAL DETAILS: Full-stack hybrid RAG benchmarking platform built with Python + Flask + Plotly Dash 2.17, deployed on HuggingFace Spaces via Docker (port 7860). No external API keys — entirely open-source. Pipeline: (1) Ingestion — arXiv API fetch by keyword, configurable batch size, 3 chunking strategies (fixed, sentence-window, semantic); (2) Embedding — all-MiniLM-L6-v2 (384-dim, L2-normalised, ~90 MB, sentence-transformers); (3) Indexing — FAISS IndexFlatIP (exact inner-product/cosine, faiss-cpu 1.8) + BM25Okapi (rank-bm25) built simultaneously; (4) Retrieval — BM25 sparse, dense vector, and hybrid all run per query; (5) Fusion — Reciprocal Rank Fusion (RRF k=60, verbatim canonical formula: 1/(60+rank)); (6) Context distillation — 6 most relevant sentences selected by sentence-level embedding similarity (not naive truncation); (7) Generation — google/flan-t5-large (780M params, seq2seq, CPU-friendly, ~900 MB, swappable); (8) Evaluation — context relevance (cosine similarity), faithfulness (answer token overlap with context), diversity (1 − mean pairwise similarity), visualised as radar + grouped bar charts. REST API endpoints: /api/fetch-arxiv, /api/query, /api/compare. All RRF, chunker, and evaluation components implemented from scratch (no LangChain). Benchmark results: hybrid RRF avg context relevance 0.47 vs 0.41 BM25; diversity 0.51 vs 0.38 BM25.
% NARRATIVE (from portfolio): Built a full-stack hybrid RAG benchmarking platform ingesting real arXiv papers on-demand, indexing with BM25 and FAISS dense retrieval simultaneously, and fusing results via Reciprocal Rank Fusion (k=60). All three retrieval methods run per query with live evaluation metrics (context relevance, faithfulness, diversity) in Plotly Dash. Flan-T5-large generates answers from distilled 6-sentence context (embedding similarity selection, not truncation). Every component — RRF, chunker, context distillation, evaluation — implemented from scratch without LangChain. Full JSON REST API for programmatic access. Hybrid RRF outperforms BM25 alone (0.47 vs 0.41 relevance; 0.51 vs 0.38 diversity). Deployed on HuggingFace Spaces via Docker, fully open-source, no API key required.
% HTML TAGS: RAG, Information Retrieval, BM25, FAISS, Dense Retrieval, Hybrid Search, Reciprocal Rank Fusion, Flan-T5, sentence-transformers, LLM Evaluation, Flask, Plotly Dash, Docker, HuggingFace, Python
\cvproject{research-rag-bench: Hybrid RAG Evaluation Platform}{https://github.com/mnoorchenar/research-rag-bench}{https://huggingface.co/spaces/mnoorchenar/research-rag-bench}{https://mnoorchenar.github.io/projects/Projects-Files/012-research-rag-bench.html}{End-to-end hybrid RAG benchmarking platform: arXiv paper ingestion → 3-strategy chunking → all-MiniLM-L6-v2 embeddings (384-dim, L2-normalised) → FAISS IndexFlatIP + BM25Okapi dual indexing → Reciprocal Rank Fusion (k=60) → 6-sentence embedding-similarity context distillation → Flan-T5-large generation (780M params); BM25, dense, and hybrid retrieval run simultaneously per query with live context relevance, faithfulness, and diversity metrics (hybrid RRF: 0.47 relevance vs 0.41 BM25); full JSON REST API; all components (RRF, chunker, evaluation) implemented from scratch without LangChain; deployed on HuggingFace Spaces via Docker, fully open-source (2025).}{\ptag{RAG}\ptag{Hybrid Search}\ptag{BM25}\ptag{FAISS}\ptag{RRF}\ptag{sentence-transformers}\ptag{Flan-T5}\ptag{LLM Evaluation}\ptag{Flask}\ptag{Plotly Dash}\ptag{Docker}}


% ── PROJECT 13 ──────────────────────────────────────────────────────────────
% DOMAIN: Agentic AI · LLM Orchestration · Retrieval-Augmented Generation · NLP · MLOps · Live Demo
% ATS KEYWORDS: agentic AI, AI agent, multi-agent system, LLM agent, autonomous agent, agent orchestration, LLM orchestration, LangGraph, LangChain, StateGraph, cyclic state machine, conditional edges, corrective RAG, CRAG, RAG, retrieval-augmented generation, hybrid search, hybrid retrieval, dense retrieval, sparse retrieval, BM25, BM25Okapi, FAISS, IndexFlatIP, cosine similarity, inner product search, Reciprocal Rank Fusion, RRF, rank fusion, sentence-transformers, BAAI/bge-small-en-v1.5, embeddings, vector embeddings, local embeddings, hallucination detection, hallucination mitigation, grounding, source citation, cited answer generation, factual accuracy, completeness evaluation, document QA, document question answering, PDF ingestion, PDF parsing, pypdf, URL scraping, BeautifulSoup, chunking, chunk overlap, Qwen, Qwen2.5-7B, open-source LLM, HuggingFace Router, OpenAI-compatible API, provider-agnostic, LCEL, LangChain Expression Language, pipe syntax, Flask, Gunicorn, Python threading, query polling, live agent trace, observability, task decomposition, research planning, LLM pipeline, prompt engineering, temperature, deterministic inference, score-based grading, keyword overlap, latency optimization, LLM call reduction, Docker, HuggingFace Spaces, MLOps, free tier deployment, NLP, natural language processing, information retrieval, knowledge retrieval, context window, top-k retrieval, top-5 recall, technical document analysis
% TECHNICAL DETAILS: Multi-agent document research system built with LangGraph 0.2 StateGraph + LangChain 0.3, deployed on HuggingFace Spaces via Docker (port 7860, free tier). Five specialised agents in linear pipeline: (1) Planner — Qwen/Qwen2.5-7B-Instruct via HF Router (OpenAI-compatible), temp 0.3, decomposes query into research plan; (2) Retriever — parallel FAISS IndexFlatIP (cosine) + BM25Okapi, fused via RRF k=60, local execution, zero API calls; (3) Grader — score-based only (hybrid score × 0.7 + keyword overlap × 0.3), 0–1 relevance, no LLM call, ~1ms latency, eliminates 5 LLM calls vs naive approach (~60% query time reduction); (4) Generator — Qwen2.5-7B, temp 0.4, max 512 tokens, structured cited answer with [Source: filename, p.N] inline citations, context-grounded only; (5) Critic — Qwen2.5-7B, temp 0.1 (near-deterministic), hallucination detection + completeness evaluation, APPROVED or NEEDS_REVIEW verdict with UI warning badge. Ingestor: PDF (pypdf, ≤10 MB) + URL (requests + BeautifulSoup); 500-token chunks, 50-token overlap; BAAI/bge-small-en-v1.5 embeddings local via sentence-transformers; FAISS + BM25 indexes built simultaneously. All LangChain chains use LCEL pipe syntax (prompt | llm). Backend: Flask 3.1 + Gunicorn + Python threading; query_id-based polling for live agent trace UI (1.5s interval). LLM layer provider-agnostic via OpenAI SDK pointed at router.huggingface.co/v1. Benchmark: hybrid RRF improves top-5 recall ~18% over pure semantic search on technical documents.
% NARRATIVE (from portfolio): Built a 5-agent LangGraph StateGraph document research system with only 3 LLM calls per query. Planner (temp 0.3) decomposes the query; Retriever runs parallel FAISS + BM25 fused via RRF (k=60); Grader scores chunks instantly via formula (no LLM, ~60% latency reduction); Generator (temp 0.4) produces cited answers grounded in retrieved context; Critic (temp 0.1) detects hallucinations and flags low-confidence answers. PDF and URL ingestion, 500-token overlapping chunks, bge-small-en-v1.5 local embeddings. Entire LLM layer provider-agnostic via OpenAI SDK; all chains use LCEL pipe syntax. Flask + threading + query-ID polling for live agent trace UI. Hybrid RRF improves top-5 recall ~18% over semantic-only. Deployed on HuggingFace Spaces via Docker, free tier.
% HTML TAGS: Agentic AI, LangGraph, LangChain, RAG, Hybrid Search, FAISS, BM25, RRF, Hallucination Detection, Qwen 2.5, LLM Orchestration, Flask, Docker, HuggingFace, Python
\cvproject{DocMind: Multi-Agent Document Research Platform with LangGraph \& Hybrid RAG}{https://github.com/mnoorchenar/docmind}{https://huggingface.co/spaces/mnoorchenar/docmind}{https://mnoorchenar.github.io/projects/Projects-Files/013-docmind.html}{LangGraph 0.2 StateGraph pipeline with 5 specialised agents (Planner, Retriever, Grader, Generator, Critic) and only 3 LLM calls per query (Qwen2.5-7B via HF Router); score-based Grader (hybrid score~×~0.7 + keyword overlap~×~0.3) eliminates 5 LLM calls ($\sim$60\% latency reduction); hybrid FAISS + BM25 + RRF retrieval (top-5 recall $+$18\% vs semantic-only); bge-small-en-v1.5 local embeddings; cited answer generation with [Source: p.N] attribution; Critic agent performs hallucination detection with APPROVED/NEEDS\_REVIEW verdict; PDF + URL ingestion; provider-agnostic LLM layer via OpenAI SDK; all chains in LCEL pipe syntax; Flask + Gunicorn + threading with live agent trace polling. Deployed on HuggingFace Spaces via Docker (2025).}{\ptag{Agentic AI}\ptag{LangGraph}\ptag{LangChain}\ptag{RAG}\ptag{Hybrid Search}\ptag{FAISS}\ptag{BM25}\ptag{RRF}\ptag{Hallucination Detection}\ptag{Qwen 2.5}\ptag{Docker}}


% ── PROJECT 14 ──────────────────────────────────────────────────────────────
% DOMAIN: Agentic AI · Conversational AI · LLM Orchestration · Customer Support Automation · MLOps · Live Demo
% ATS KEYWORDS: agentic AI, AI agent, ReAct, ReAct agent, Reason and Act, LangGraph, LangChain, StateGraph, conditional edges, multi-turn, multi-turn chatbot, conversational AI, customer support automation, customer support chatbot, chatbot, virtual assistant, intelligent agent, LLM orchestration, tool use, tool calling, function calling, tool dispatch, agent tools, agentic workflow, agent loop, iteration depth, FAQ search, order tracking, ticket creation, knowledge base, escalation, product search, HuggingFace Inference API, InferenceClient, streaming, token streaming, Server-Sent Events, SSE, real-time streaming, live streaming, Mistral, Zephyr, Phi-3, Llama 3, open-source LLM, free LLM, multi-model, runtime model switching, regex parsing, ReAct parsing, structured output parsing, JSON fallback, TypedDict, AgentState, Flask, gunicorn, gthread, Python threading, daemon thread, queue.Queue, thread-safe, non-blocking SSE, session state, in-memory state, zero-persistence, Docker, HuggingFace Spaces, Chart.js, session analytics, graph trace, observability, node timing, latency tracking, tool usage analytics, response latency, token estimation, Python, NLP, natural language processing, LLM deployment, MLOps, production chatbot, temperature, low variance inference
% TECHNICAL DETAILS: Multi-turn ReAct customer support chatbot built with LangGraph 0.2 StateGraph + LangChain Core, deployed on HuggingFace Spaces via Docker (port 7860, free tier). Four-node pipeline: (1) Router — initialises state, resets iteration counters; (2) Agent — HuggingFace Hub InferenceClient chat_completion(stream=True), temp 0.25, regex-based ReAct parser (Action/Action Input blocks; JSON fallback for malformed output); (3) Tool Executor — dispatches one of 5 tools (search_faq: keyword-scored KB, 15 entries; check_order_status; create_support_ticket; search_products: 8 products; escalate_to_human); conditional edge loops back to Agent for up to 4 iterations, then routes to Responder; (4) Responder — emits done SSE event and session analytics. LLMs: 4 free HF models selectable at runtime mid-session (Mistral-7B-Instruct, Zephyr-7B-β, Phi-3-Mini-4k-Instruct, Meta-Llama-3-8B-Instruct). SSE architecture: LangGraph runs in daemon thread; per-session queue.Queue bridges to Flask SSE generator; each node emits enter/exit events with time.time() timing (±1ms). UI: 6 modules — streaming chat, live animated graph trace (4 nodes, pending/running/completed states), collapsible tool call log (name + input JSON + output + latency), Chart.js session analytics (horizontal bar: tool frequency; line: latency history), conversation history (role badge, 280-char preview, token estimate), multi-model selector. Backend: Flask 3 + gunicorn 1 worker × 4 gthread. Token count: word count × 1.35 estimate. Zero-persistence: all state in Python dict, no DB, no file writes.
% NARRATIVE (from portfolio): Built a multi-turn LangGraph 0.2 ReAct customer support agent with 4-node StateGraph (Router → Agent → Tool Executor → Responder), conditional looping up to 4 iterations, and 5 live tools (FAQ search, order lookup, ticket creation, product search, human escalation). Free-tier HuggingFace models (Mistral, Zephyr, Phi-3, Llama 3) selectable at runtime; regex-based ReAct parser with JSON fallback handles unstructured open-model output. SSE token streaming via per-session queue.Queue bridging LangGraph daemon thread to Flask response. Live animated graph trace with per-node ms timing. Chart.js session analytics (tool frequency + latency history). Zero-persistence, in-memory architecture for indefinite free-tier deployment. Deployed on HuggingFace Spaces via Docker.
% HTML TAGS: Agentic AI, LangGraph, ReAct, LangChain, Conversational AI, Chatbot, SSE Streaming, Tool Use, LLM Orchestration, HuggingFace, Flask, Docker, Customer Support AI, Python, Multi-Turn
\cvproject{TechStore Support Agent: Multi-Turn ReAct Chatbot with LangGraph}{https://github.com/mnoorchenar/langgraph-support-agent}{https://huggingface.co/spaces/mnoorchenar/langgraph-support-agent}{https://mnoorchenar.github.io/projects/Projects-Files/014-langgraph-support-agent.html}{LangGraph 0.2 ReAct StateGraph (4 nodes: Router → Agent → Tool Executor → Responder) with conditional looping up to 4 iterations and 5 live tools (FAQ search, order lookup, ticket creation, product search, human escalation); 4 free HuggingFace LLMs switchable at runtime (Mistral-7B, Zephyr-7B, Phi-3-Mini, Llama-3-8B); regex-based ReAct parser with JSON fallback for open-model output; SSE token streaming via per-session \texttt{queue.Queue} bridging LangGraph daemon thread to Flask; live animated 4-node graph trace with per-node ms timing; Chart.js session analytics (tool frequency + response latency); zero-persistence in-memory architecture. Deployed on HuggingFace Spaces via Docker (2025).}{\ptag{Agentic AI}\ptag{LangGraph}\ptag{ReAct}\ptag{LangChain}\ptag{Tool Use}\ptag{SSE Streaming}\ptag{Conversational AI}\ptag{HuggingFace}\ptag{Flask}\ptag{Docker}\ptag{Python}}